\notechapter{1}{Operator spaces, compression, and alignment}

\section*{Scope}

The replacement of a first-principles Hamiltonian by a finite model is not
defined solely by retaining a selected set of eigenvalues. A reduced
description must specify the state space on which the reduced operator acts,
the map from the parent space to that state space, the coordinate
representation used for its matrix elements, and the class of operators
admitted as candidate models. These choices determine which algebraic
operations are meaningful and which physical interpretations can be assigned
to their results.

The distinction is particularly important when the parent Hamiltonian is an
unbounded differential operator, while the reduced description is a finite
matrix. In the continuum, the operator domain and boundary conditions are part
of the Hamiltonian. In the retained representation, those data are no longer
written as an explicit domain, but their consequences remain encoded in the
matrix. A rigorous reduction must therefore distinguish the continuum
operator, the retained operator, and the numerical matrix representing that
operator.

\section*{Operators and their domains}

Let $\mathcal H$ be a complex Hilbert space. An unbounded operator is a map
\begin{equation}
    \hat A:
    \mathcal D(\hat A)\subseteq\mathcal H
    \longrightarrow
    \mathcal H,
\end{equation}
where the domain $\mathcal D(\hat A)$ is part of the definition of $\hat A$.
Two operators may have the same formal differential action and nevertheless
define different physical systems because they act on different domains. For
quantum Hamiltonians, the distinction between a symmetric differential
expression and a self-adjoint operator is essential: self-adjointness requires
equality of both the operator action and the adjoint domain~\cite{teschl2014,reedSimon1975,schmudgen2012}.

The sum or difference of two unbounded operators is initially defined on the
intersection of their domains,
\begin{equation}
    \mathcal D(\hat A\pm\hat B)
    =
    \mathcal D(\hat A)\cap\mathcal D(\hat B).
\end{equation}
Consequently, an expression such as
\begin{equation}
    \hat V=\hat H-\hat T
\end{equation}
is an operator identity only after the operators and their compatible domains
have been specified. This qualification is not an objection to subtraction.
It identifies the mathematical data required for the subtraction to refer to
a common operator space.

\section*{Spectral restriction}

Consider a self-adjoint Hamiltonian $\hat H$ with a pure-point spectral sector,
\begin{equation}
    \hat H\lvert\psi_i\rangle
    =
    E_i\lvert\psi_i\rangle,
\end{equation}
where the retained eigenvectors are orthonormal. Let $\mathcal I_P$ denote the
selected set of spectral indices and define
\begin{equation}
    \hat P
    =
    \sum_{i\in\mathcal I_P}
    \lvert\psi_i\rangle\langle\psi_i\rvert,
\end{equation}
with retained space
\begin{equation}
    \mathcal H^{(P)}
    =
    \operatorname{Ran}\hat P.
\end{equation}
If an energy cutoff intersects a degenerate eigenspace, the full degenerate
eigenspace must be retained if the construction is to remain invariant under
unitary rotations within that eigenspace. More generally, the selected range
should be an invariant subspace of the parent operator.

The retained Hamiltonian is
\begin{equation}
    \hat H^{(P)}
    =
    \left.
        \hat P\hat H\hat P
    \right|_{\mathcal H^{(P)}}:
    \mathcal H^{(P)}
    \longrightarrow
    \mathcal H^{(P)}.
\end{equation}
It reproduces the action of the parent Hamiltonian exactly on the selected
invariant subspace. When $\hat P\hat H\hat P$ is instead regarded as an
operator on the full parent space, it acts as zero on
$\operatorname{Ran}(\hat I-\hat P)$. That zero action is an embedding
convention; it is not the parent Hamiltonian on the discarded sector.

This distinction separates three objects:
\begin{equation}
    \hat H,
    \qquad
    \hat H^{(P)},
    \qquad
    \mathbf H^{(P)}.
\end{equation}
They denote, respectively, the parent operator, the retained operator, and a
matrix representation of the retained operator in a declared basis. Treating
them as interchangeable obscures the state-space maps on which later
comparisons depend.

\section*{Compression and downfolding}

A fixed-subspace compression should not be identified with an effective
Hamiltonian obtained by eliminating the complementary sector. With
\begin{equation}
    \hat Q=\hat I-\hat P,
\end{equation}
a Feshbach--L\"owdin construction gives the energy-dependent operator
\begin{equation}
    \hat H_{\mathrm{down}}(E)
    =
    \hat P\hat H\hat P
    +
    \hat P\hat H\hat Q
    \left(E-\hat Q\hat H\hat Q\right)^{-1}
    \hat Q\hat H\hat P,
\end{equation}
whenever the reduced resolvent exists~\cite{lowdin1982}. The second term
represents virtual excursions into the eliminated subspace. In general,
\begin{equation}
    \hat H_{\mathrm{down}}(E)
    \neq
    \hat H^{(P)}.
\end{equation}
The approximation target must therefore be declared explicitly. A model
fitted to $\hat H^{(P)}$ approximates the parent action on a selected subspace.
A model fitted to $\hat H_{\mathrm{down}}(E)$ also attempts to reproduce
eliminated-space effects and may inherit energy dependence and nonlocality.

\section*{Common compression and additive structure}

Suppose that
\begin{equation}
    \hat H=\hat T+\hat V
\end{equation}
is valid on a compatible domain, and assume that the finite-dimensional
retained space lies within the required domains. Applying the same compression
to every term gives
\begin{equation}
    \hat H^{(P)}
    =
    \hat T^{(P)}+\hat V^{(P)},
\end{equation}
where
\begin{equation}
    \hat T^{(P)}
    =
    \left.\hat P\hat T\hat P\right|_{\mathcal H^{(P)}},
    \qquad
    \hat V^{(P)}
    =
    \left.\hat P\hat V\hat P\right|_{\mathcal H^{(P)}}.
\end{equation}
It follows that
\begin{equation}
    \hat V^{(P)}
    =
    \hat H^{(P)}-\hat T^{(P)}.
\end{equation}
The subtraction is valid because the corresponding addition was performed in
the common vector space
\begin{equation}
    \mathcal L\!\left(\mathcal H^{(P)}\right).
\end{equation}

The physical interpretation of the difference is a separate question.
Although $\hat V$ may be a local multiplication operator in the parent
position representation, $\hat P\hat V\hat P$ is generally nonlocal. If
\begin{equation}
    \langle x\rvert\hat V\lvert x'\rangle
    =
    V(x)\delta(x-x'),
\end{equation}
then
\begin{equation}
    V^{(P)}(x,x')
    =
    \int dy\,
    P(x,y)V(y)P(y,x'),
\end{equation}
which is not generally proportional to $\delta(x-x')$. Locality must therefore
be treated as an admissible model assumption rather than as an automatic
consequence of projection.

\section*{Hilbert--Schmidt geometry}

The space of all bounded operators on an infinite-dimensional Hilbert space is
not a Hilbert space under the trace inner product. The relevant
infinite-dimensional space is the Hilbert--Schmidt class
\begin{equation}
    \mathfrak S_2(\mathcal H)
    =
    \left\{
        \hat A\in\mathcal B(\mathcal H)
        \;\middle|\;
        \operatorname{Tr}(\hat A^\dagger\hat A)<\infty
    \right\},
\end{equation}
with
\begin{equation}
    \langle\hat A,\hat B\rangle_{\mathrm{HS}}
    =
    \operatorname{Tr}(\hat A^\dagger\hat B)
\end{equation}
and
\begin{equation}
    \lVert\hat A\rVert_{\mathrm{HS}}^2
    =
    \operatorname{Tr}(\hat A^\dagger\hat A)
\end{equation}
\cite{simon2005}. Full continuum Hamiltonians are generally unbounded and are
not Hilbert--Schmidt operators. The trace inner product should not be applied
to them without an additional construction.

Finite-dimensional retained operators avoid this obstruction. In any
orthonormal basis of an $M$-dimensional retained space,
\begin{equation}
    \lVert\hat A\rVert_{\mathrm{HS}}^2
    =
    \sum_{\alpha,\beta=1}^{M}
    \lvert A_{\alpha\beta}\rvert^2
    =
    \lVert\mathbf A\rVert_{\mathrm F}^2.
\end{equation}
The Frobenius norm is therefore the coordinate representation of the
Hilbert--Schmidt norm. Its unitary invariance makes it suitable for operator
comparisons after the relevant state spaces have been identified. It does not
determine the identification itself.

\section*{Alignment across retained spaces}

Let $\hat H_b^{(P)}$ and $\hat H_d^{(P)}$ denote retained pristine and doped
Hamiltonians acting on $\mathcal H_b^{(P)}$ and $\mathcal H_d^{(P)}$. Equal
matrix dimensions do not imply that the two matrices represent operators on
the same physical state space. A unitary identification
\begin{equation}
    \hat U_d:
    \mathcal H_b^{(P)}
    \longrightarrow
    \mathcal H_d^{(P)}
\end{equation}
must first be specified. The aligned difference is
\begin{equation}
    \Delta\hat H_d^{(P)}
    =
    \hat U_d^\dagger
    \hat H_d^{(P)}
    \hat U_d
    -
    \hat H_b^{(P)}.
\end{equation}

Wannier functions provide localized frames for such comparisons, but their
construction includes unitary gauge freedom, and an entangled band manifold
requires an explicit retained-subspace selection~\cite{marzari2012,souza2001}.
The alignment record must therefore include the retained subspaces, orbital
labels, gauge, geometry, units, spin convention, and energy reference.

The fundamental conclusion is that model reduction is performed on identified
operator spaces, not on arrays whose dimensions happen to agree. Once the
spaces and maps are fixed, subtraction and Hilbert--Schmidt comparison are
mathematically straightforward. The scientific content lies in justifying the
projection, the identification, and the admissible model class.
